\section{Repository integration and artifact guide}
\label{sec:integration}

\subsection{Compatibility and public interfaces}

The archive contains a self-contained copy of the pinned \texttt{fast/}
package, an exact baseline copy, and a patch with paths relative to the
ProveIt repository root. The implementation needs Python 3.10 or newer
and has no new third-party runtime dependency. The recorded checks use
Python 3.12.14. Matplotlib is needed only to regenerate the research
figure; the figure and generated tables are already included for a
standalone PDF build.

The public scanner options are
\begin{lstlisting}[language=Python]
FastScan(..., composition="legacy")
khovanov_rank(..., composition="legacy")
recognize(..., composition="legacy",
          quotient_max_assignments=0,
          quotient_seconds=0.1)
\end{lstlisting}
The accepted composition modes are \texttt{legacy}, \texttt{adaptive},
and \texttt{dense}. The latter forces the new dense paths while retaining
exact zero, identity, and self-product shortcuts. Nonlegacy modes require
the bits algebra, min-fill cancellation, and self-inverse option where
that option is exposed. Unsupported combinations raise an error. The
implementation does not silently ignore them. Competing worker processes
and factored calls preserve the chosen composition mode, shape-cache
setting, and optional $d^2$ checks.

The finite-target filter is disabled when its assignment budget is zero.
A positive budget applies to complete seed-orbit assignments per examined
summand, across the configured palette. The default cooperative filter
cap is 0.1 seconds; the Python API permits \texttt{None} to remove that
separate cap. A global recognition deadline still limits the filter by
the remaining time. Exhaustion of the filter's own budget records
\texttt{INCONCLUSIVE} and continues. Exhaustion of the global recognition
budget returns \texttt{UNKNOWN}. These distinctions are part of the
public behavior, not merely diagnostic logging.

Example commands, run from \texttt{implementation/fast/}, are
\begin{lstlisting}[language=bash]
python -m fastunknot khovanov examples/conway.json \
  --composition adaptive --check-d2

python -m fastunknot recognize examples/conway.json \
  --no-jones --quotient-max-assignments 5000 \
  --quotient-seconds 0.1

python -m unittest discover -s tests -v
\end{lstlisting}
The second command disables Jones specifically to expose the new
certificate stage on a fixture that the default Jones filter already
handles. It is a demonstration of the fallback, not a recommended change
to filter order.

\subsection{Witness evidence and independent replay}

A successful finite-target search returns method
\texttt{finite-quotient-A5-seeds}. Evidence contains the exact normalized
working \texttt{diagram\_pd}, certificate, checked seed plan, orbit counts,
attempt details, and status. When preprocessing factors a diagram,
the witness belongs to a specified summand and its evidence is stored
with that summand. Existing simplification and factorization correctness
connect that working diagram to the original input. The finite-target
checker itself verifies precisely the PD supplied to it; it does not
claim to replay the earlier Reidemeister or connected-sum transformations.

For an unfactored result whose finite-target stage succeeded, replay is
\begin{lstlisting}[language=Python]
from fastunknot.finite_quotient_check import verify_certificate

evidence = result.evidence["finite_quotient"]
checked = verify_certificate(evidence["diagram_pd"],
                             evidence["certificate"])
assert checked["valid"]
\end{lstlisting}
The supplement also provides a portable helper that replays every
archived certificate against its exact fixture. Keeping explicit
permutations and the precise PD is essential: an image-group order or
an opaque solver state is insufficient evidence by itself.

\subsection{What each directory contains}

\begin{center}
\small
\begin{tabular}{P{.35\linewidth}P{.57\linewidth}}
\toprule
Path & Purpose \\
\midrule
\texttt{article/} & Main TeX file, section sources, generated tables,
vector figure, bibliography, and compiled PDF. \\
\texttt{implementation/fast/} & Runnable package, optional backends,
74-test suite, isolated scanner benchmark, and integration notes. \\
\texttt{baseline/fast/} & Forty source files verified against their
pinned Git blob identities; prior result directories are omitted. \\
\texttt{verification/algebra/} & Portable dense benchmark, scalar
comparison suite, associativity/transfer checks, grading audit, and raw
measurement JSON. \\
\texttt{verification/finite\_quotient/} & Search comparison, certificates,
seed plans, exact fixture bytes, exhaustive negative controls, and replay
helper. \\
\texttt{provenance/} & Repository revisions, baseline blob records,
source normalization details, and validation records. \\
\texttt{tools/} & Table and figure generation, archive verification,
and patch reproduction helpers. \\
\texttt{patches/} & Integration patch against the pinned repository
paths. \\
\bottomrule
\end{tabular}
\end{center}

The top-level README gives exact commands from a freshly extracted
archive. The TeX bibliography is self-contained; no BibTeX or Biber
database download is required. Compiling the main file with ordinary
pdfLaTeX passes resolves the table of contents and cross-references.
The SHA-256 manifest covers the delivered source, data, and PDF files.
The verification helper checks the manifest, the original Git blob
identities, recorded fixture hashes, and the correspondence of the patch
to the supplied baseline and implementation.

\subsection{Merge boundaries}

The intentional edits to existing implementation files are limited to
the scanner selector, public option forwarding and counters, optional
recognizer stage and shared budget, CLI propagation, and the README link.
New modules implement dense composition, finite-target search, and its
separate verifier. The patch includes the new tests and integration
notes. Research measurements and the article may be stored under a new
report directory when integrating the archive.

The patch is concrete and reviewable against the pinned revision. If the
repository has advanced, its changes should be reconciled with the newer
code and the same suite rerun. In particular, report-07 sharing, the
pointed scanner, three-braid specialization, and newer preprocessing
branches are not represented here as an already tested combined release.
The theorem about graded component types applies to the specified exact
sharing construction; implementation and performance of that combined
backend remain a separate, well-defined task.
